\begin{abstract}
Este trabalho investiga um problema de Pesquisa Operacional motivado por
operações logísticas em que as unidades de carga estão dispostas em uma
sequência física de acesso. Nessas situações, o operador deve decidir como
subsequências contíguas de cargas devem ser atribuídas a uma frota heterogênea,
respeitando restrições operacionais e econômicas. Embora a motivação inicial
tenha surgido no transporte de bobinas de aço, a estrutura do problema é mais
geral. Ela se aplica a contextos em que a atribuição entre carga e veículo fixa
o carregamento, sem gerar um subproblema adicional de arranjo interno. Como
etapa de posicionamento teórico, o estudo inclui uma revisão rápida da
literatura sobre alocação de frota, atribuição de veículos e planejamento de
transporte de cargas, a partir da qual se delimita o espaço da contribuição do
modelo em desenvolvimento. Em seguida, o problema é descrito por uma formulação
inteira do tipo \emph{set partitioning over intervals}, por uma versão
simplificada em programação dinâmica e por sua extensão completa com controle
explícito da frota, de modo a explicitar a estrutura sequencial da decisão e o
papel das restrições globais de recursos. Como desdobramento aplicado, o estudo
também inclui a implementação de um protótipo interativo da formulação
completa, usado como base para testes em campo e para a coleta estruturada de
instâncias e soluções geradas durante o uso do algoritmo.
\end{abstract}
